%% kapa: exam=1 round=21 year=2010 period=1 form=A subject=economics subject_ko=경제원론 seq=14 number=54
%% source: Q-Net 감정평가사 기출문제 (cst00309) artlSeq=5024181 file=extracted/1008/5024181/2010년 제21회 감정평가사 1차시험 기출문제/1교시.hwp
%% conversion: kordoc 4.24.51 -> tools/gen1.py (auto); review=proofread-sonnet; flags=none
\begin{question}{54}
불완전경쟁시장에 관한 설명으로 옳지 않은 것은? (단, 개별 기업의 생산량은 0보다 크다.)
\begin{choices}
  \item 독점기업의 한계수입(MR)과 평균수입(AR)은 MR < AR의 관계를 갖는다.
  \item 독점적 경쟁기업의 장기균형에서는 평균수입과 장기평균비용이 일치한다.
  \item 과점시장의 모든 기업들이 담합해 카르텔을 결성하면 실질적으로 독점체제를 구축하는 결과를 가져온다.
  \item 독점적 경쟁기업은 장기적으로 ‘가격 = 한계비용’인 생산량 수준에서 이윤을 극대화한다.
  \item 독점기업이 가격차별을 하게 되면 가격탄력성이 높은 소비자는 가격탄력성이 낮은 소비자보다 더 낮은 가격을 지불하게 된다.
\end{choices}
\end{question}
\answer{4}
